\documentclass[12pt]{article}

\usepackage{geometry}
\usepackage{longtable}
\usepackage{hyperref}
\usepackage{array}

\geometry{margin=1in}

\title{
    \textbf{Snapshot Objectives Overview} \\
    JPL Lunar Detection Pipeline
}
\author{
    Software Engineering Tools Final Project \\
    Course: CS3338 \\
    Team Members: [Name 1], [Name 2], [Name 3], [Name 4]
}
\date{[Insert Date]}

\begin{document}

\maketitle
\newpage

\tableofcontents
\newpage

\section{Introduction}

This document gives an overview of the four project snapshots for the JPL Lunar Detection Pipeline. Each snapshot represents a stage of the simulated software engineering project.

The project is a simulated open-source computer vision toolset for detecting lunar surface features such as craters, boulders, and geological structures. The purpose of the snapshots is to show how the project grows over time through planning, documentation, mock data, Jira tasks, Docker planning, and TestRail testing.

\section{Snapshot Summary}

\begin{longtable}{|p{1.5in}|p{2in}|p{2.5in}|}
\hline
\textbf{Snapshot} & \textbf{Main Goal} & \textbf{Main Deliverables} \\
\hline
Snapshot 1 & Project foundation & GitHub repo, SDD, SRS, README, Design Spec, User Manual, initial mock data \\
\hline
Snapshot 2 & Add boulder detection feature & Updated documents, Jira tasks, TestRail report, boulder detection mock data \\
\hline
Snapshot 3 & Add crater filtering and history feature & Updated documents, Jira tasks, TestRail report, crater filtering, detection history \\
\hline
Snapshot 4 & Finalize project and future work & Final documents, final TestRail report, project reflection, future work \\
\hline
\end{longtable}

\section{Snapshot Timeline}

The project will be completed in four stages:

\begin{enumerate}
    \item \textbf{Snapshot 1:} Create the project foundation.
    \item \textbf{Snapshot 2:} Add the first major feature and begin TestRail documentation.
    \item \textbf{Snapshot 3:} Add a second major feature and expand testing.
    \item \textbf{Snapshot 4:} Finalize all project documents and reflect on future work.
\end{enumerate}

\section{Tools Used Across Snapshots}

The project will use the following tools:

\begin{itemize}
    \item \textbf{GitHub:} Stores the repository and project documents.
    \item \textbf{Docker:} Describes expected system services using Docker Compose.
    \item \textbf{Jira:} Tracks sprint tasks, bugs, and project progress.
    \item \textbf{TestRail:} Stores test cases and test run reports.
    \item \textbf{LaTeX:} Used to write formal project documents.
\end{itemize}

\section{Conclusion}

The snapshot objective documents help explain how the project changes over time. Each snapshot adds new work, updates existing documents, and shows progress toward the final project submission.

\end{document}